\section{Related work and claim boundary}

The workflow combines established research practices rather than claiming them individually. Preregistration formalizes the separation between hypotheses/analyses defined before outcomes and post-outcome explanations \cite{nosek2018preregistration}. Scientific-computing guidance emphasizes version control, testing, automation, and recoverable workflows \cite{wilson2014best}; data-provenance research formalizes how derived outputs remain traceable to source records \cite{buneman2001provenance}. Tool-using autonomous scientific systems have also demonstrated that language-model agents can plan and execute multi-step research activities \cite{boiko2023autonomous}. These literatures motivate components of the present work but do not establish the particular recursive negative-recovery contract studied here.

The candidate residual is their executable composition around recursive negative-result recovery: a typed terminal system in which donor absorption, refutation, support, and cannot-check remain distinct; reopening is tied to a recorded residual; and the evidence path is receipt-bound. That residual is conditional on a fresh submission-date literature search. The manuscript therefore avoids `first', `unique', or exhaustive-nearest-work language until that search is complete.

The exact TARE expressivity result, the controller--host agreement benchmark, and the typed/scoped partial-knowledge mechanism evidence are each owned by separate reports and are cited here as ordinary references. Systems-level guarantees about the execution harnesses are a further subject, and are claimed by none of them. This paper cites those outcomes only to show the behaviour of the common methodology.